\section{A note on Mathlib}\label{sec:00-setup:04-mathlib-note:1}

This book builds groups, rings, and path algebras \textbf{from scratch}, deliberately, without importing Mathlib (Lean's community math library). This is slower, but better for learning: every definition and every proof obligation is made explicit. Chapter 13 points toward Mathlib for readers ready to use the ``real'' library instead of reinventing it.

Starting in Chapter 6, most worked examples are followed by a small, clearly labeled ``Mathlib equivalent'' box, showing the same statement written against Mathlib's actual \texttt{Group}/\texttt{Ring}/\texttt{Module} API. This does not replace the from-scratch approach; the hand-built version remains the main teaching path, and the Mathlib box is only a preview. Holding both versions in mind at once --- the definition just derived, and the shape the same idea takes in the library used later --- builds a sharper understanding of both than either alone would.

\textbf{Key points.} Lean 4 plus \texttt{lake}, an editor with the Lean extension, and a pinned toolchain (matching \texttt{lean\_\allowbreak{}project/\allowbreak{}lean-\allowbreak{}toolchain}) are all that is needed to follow along. This book is Mathlib-free by design through Chapter 11's from-scratch constructions; Mathlib appears only in the ``Mathlib equivalent'' boxes from Chapter 6 onward, and in full starting Chapter 13.

\textbf{Socratic questions.}

\begin{enumerate}
\def\labelenumi{\arabic{enumi}.}
\tightlist
\item
  \emph{\texttt{lean\_\allowbreak{}project} already has Mathlib installed as a dependency (it is what powers the ``Mathlib equivalent'' boxes from Chapter 6 onward) --- so why not just import it everywhere from page one, instead of building \texttt{Group}/\texttt{Ring} from scratch first?} Because the point of this book is not merely to \emph{use} a group in Lean, but to see exactly what a group \emph{is} to Lean: every field, every proof obligation, with nothing hidden behind someone else's typeclass hierarchy. A library saves effort precisely by hiding that machinery --- which is exactly the wrong thing to hide from a first encounter.
\item
  \emph{\texttt{elan} pins one exact Lean version per project via \texttt{lean-\allowbreak{}toolchain}. What would go wrong without that pin, on a machine with several Lean projects at once?} A later toolchain update to one project could silently change how another project's code elaborates or even fails to compile --- the entire point of \texttt{lean\_\allowbreak{}project/\allowbreak{}lean-\allowbreak{}toolchain} reading \texttt{leanprover/\allowbreak{}lean4:v4.31.0} is that every code block in this book stays reproducible regardless of what else is installed system-wide.
\item
  \emph{If Mathlib is Mathlib-free by design through Chapter 11, why does Chapter 6 onward show Mathlib code at all?} Because ``from scratch'' and ``never see the real library'' are different promises --- this book keeps the first and deliberately breaks the second, by pairing every hand-built definition with a labeled preview of its Mathlib counterpart, so the transition in Chapter 13 is a recognition, not a cold start.
\end{enumerate}

